\section{Workflow Summary}

\begin{table}[H]
\centering
\small
\begin{tabular}{lll}
\toprule
\wcell{3.2cm}{\textbf{Practice}} & \wcell{4.3cm}{\textbf{How it was applied}} & \wcell{4.1cm}{\textbf{Why it mattered}} \\
\midrule
\wcell{3.2cm}{Canonical remote repository} & \wcell{4.3cm}{The project was consolidated into one remote repository with visible branch and merge history.} & \wcell{4.1cm}{Provided a dated engineering trail rather than isolated local files.} \\
\wcell{3.2cm}{Branch-based change integration} & \wcell{4.3cm}{Major report and code changes were grouped into identifiable commits and merge checkpoints.} & \wcell{4.1cm}{Reduced ambiguity about what changed and when.} \\
\wcell{3.2cm}{Review-first code changes} & \wcell{4.3cm}{Later submission-facing changes were aligned to a branch-and-review workflow.} & \wcell{4.1cm}{Improved accountability for risky integration work.} \\
\wcell{3.2cm}{Evidence-tree preservation} & \wcell{4.3cm}{Testing records, issue logs, screenshots, and report material were kept in the documentation tree.} & \wcell{4.1cm}{Made the report easier to support with named artifacts.} \\
\bottomrule
\end{tabular}
\caption{Version-control and repository workflow used by the project.}
\label{tab:appendix-version-control}
\end{table}

\section{Workflow Limits}

The repository trail is stronger than the earliest coordination trail. Early planning relied more on informal coordination than on a fully preserved digital board history. The final report therefore distinguishes between:

\begin{itemize}
    \item repository activity that is directly provable from commits and merges;
    \item planning or meeting context that is preserved only partially;
    \item recovered notes that remain labelled as recovered.
\end{itemize}

This appendix keeps that distinction visible instead of implying stronger historical evidence than the workspace can prove.
